\begin{sol}{xca:center}
    Clearly $1\in Z(G)$. If $x\in Z(G)$, then $xg=gx$ for all $g\in G$. Multiplying by $x^{-1}$ on the left and 
    on the right, we get that $gx^{-1}=x^{-1}$ holds for all $g\in G$. Finally, if $x,y\in Z(G)$. Then
    \[
    (xy)g=x(yg)=x(gy)=(xg)y=(gx)y=g(xy)
    \]
    for all $g\in G$. Hence $xy\in Z(G)$. 
\end{sol}

% \begin{sol}{xca:centralizer}
%     First, $1\in G_G(g)$, as $1g=g1=g$. If $x\in c_G(g)$, then $xg=gx$. Multiplying 
%     on the left and on the right by $x^{-1}$, one gets $x^{-1}g=gx^{-1}$, that is 
%     $x^{-1}\in C_G(g)$. Finally, if $x,y\in C_G(g)$, then 
%     \[
%     $(xy)g=$
%     \]
% \end{sol}

\begin{sol}{xca:conjugate}
    Since $S$ is a subgroup, $1\in S$ and 
    if $x,y\in S$, then $x^{-1}\in S$ and $xy\in S$. Now 
    $1\in gSg^{-1}$, as $1\in S$ and $1=g1g^{-1}$. If $x\in gSg^{-1}$, then 
    $x=gsg^{-1}$ for some $s\in S$. Thus 
    \[
    x^{-1}=(gsg^{-1})^{-1}=gs^{-1}g^{-1}\in gSg^{-1},
    \]
    as $s^{-1}\in S$. Finally, 
    if $x=gsg^{-1}\in gSg^{-1}$ and $y=gtg^{-1}\in gSg^{-1}$ for some $s,t\in S$, then 
    \[
    xy=(gsg^{-1})(gtg^{-1})=g(st)g^{-1}\in gSg^{-1},
    \]
    as $st\in S$. 
\end{sol}

\begin{sol}{xca:center_S3}
    If $\sigma\in Z(\Sym_3)$ and $\sigma\ne\id$, there exists $i\in\{1,2,3\}$ such that 
    $\sigma(i)\ne i$. Let $j=\sigma(i)$ and $k\in\{1,2,3\}\setminus\{i,j\}$. Then 
    $(jk)\sigma$ is a permutation such that $i\mapsto k$, while 
    $\sigma(jk)$ is such that $i\mapsto j$. In particular, $(jk)\sigma\ne\sigma(jk)$, a contradiction.

    The group $\Sym_3$ has six elements: $\id$, $(12)$, $(13)$, $(23)$. $(123)$ and $(132)$. 
    First note that $\id\in C_{\Sym_3}((12))$ and 
    $(12)\in C_{\Sym_3}((12))$. However, 
    the permutations $(23)$, $(13)$, $(123)$ and $(132)$ do not commute with
    $(12)$. For example, 
    \[
    (23)(12)=(132)\ne (123)=(12)(23).
    \]
\end{sol}

\begin{sol}{xca:subgroup}
    Let us prove $\implies$. Since $1\in S$, then $S\ne\emptyset$. If $u,v\in S$, then 
    $v^{-1}\in S$ and $uv^{-1}\in S$. 

    Let us prove now $\impliedby$. If $S\ne\emptyset$, let $u\in S$. Then $1=uu^{-1}\in S$. The assumption 
    Let $u,v\in S$. The assumption with $x=1\in S$ and $y=v$ yields $v^{-1}\in S$. The assumption 
    with $x=u$ and $y=v^{-1}$ yields $uv\in S$. 
\end{sol}

\begin{sol}{xca:SL_subgroup}
    The identity matrix belongs to $\SL_n(\R)$. If $a,b\in\SL_n(\R)$, then 
    $ab^{-1}\in\SL_n(\R)$, as 
    \[
    \det(ab^{-1})=\det(a)\det(b^{-1})=\det(a)\det(b)^{-1}=1.
    \]
    By Exercise \ref{xca:subgroup}, $\SL_n(\R)$ is a subgroup of $\GL_n(\R)$. 
\end{sol}

\begin{sol}{xca:intersection}
    Let $\{H_\lambda:\lambda\in\Lambda\}$ be a collection of subgroups of a group $G$ and 
    $H=\cap_{\lambda\in \Lambda}H_\lambda$. We claim that $H$ is a subgroup of $G$. Since
    $1\in H_\lambda$ for all $\lambda$, $H$ is non-empty. If $x,y\in H$, then $x,y\in H_\lambda$ for all $\lambda$. 
    Since each $H_\lambda$ is a subgroup of $G$, $xy^{-1}\in H_\lambda$ for all $\lambda$. Thus $xy^{-1}\in H$.
\end{sol}

\begin{sol}{xca:generated}
    Let 
    \[
    H=\{x_1^{n_1}\cdots x_k^{n_k}:k\geq0,\,x_1,\dots,x_k\in X,\,-1\leq n_1,\dots,n_k\leq 1\}.
    \]
    To prove that $H\subseteq\langle X\rangle$, let $h=x_1^{n_1}\cdots x_k^{n_k}\in H$. 
    If $S$ is a subgroup of $G$ containing $X$, then $x_j\in S$ for all $j$. This implies that 
    $h=x_1^{n_1}\cdots x_k^{n_k}\in S$. Thus 
    \[
    h\in\bigcap_{\substack{S\leq G\\X\subseteq S}}S.
    \]
    
    To prove that $H\supseteq \langle X\rangle$ we first 
    claim that $H$ is a subgroup of $G$. Note that $H\ne\emptyset$, as $1\in H$ (this is the empty word). If 
    $u=x_1^{n_1}\cdots x_k^{n_k}\in H$ and 
    $v=x_{k+1}^{n_{k+1}}\cdots x_{l}^{n_l}\in H$, then 
    \[
    uv^{-1}=x_1^{n_1}\cdots x_k^{n_k}x_{l}^{-n_{l}}\cdots x_{k+1}^{-n_{k+1}}\in H. 
    \]
    Now note that $H$ is a subgroup of $G$ containing $X$. Thus 
    \[
    \langle X\rangle=\bigcap_{\substack{S\leq G\\X\subseteq S}}S\subseteq H.
    \]
\end{sol}

\begin{sol}{xca:union}
    Let $G=\Sym_3$. Then $H=\{\id,(12)\}$ and 
    $K=\{\id,(23)\}$ are subgroups of $G$. However, 
    $H\cup K=\{\id,(12),(23)\}$ is not a subgroup, as 
    $(12)(23)=(123)\not\in H\cup K$. 
\end{sol}


